\documentclass{article}
\usepackage[utf8]{inputenc}
\usepackage{a4wide}

\begin{document}

\section*{เตรียมระดมทุน}

สมมติว่า บริษัทซุปเปอร์ต๋อยจะเริ่มเสนอขายหุ้นในการจดทะเบียนเข้าตลาดหลักทรัพย์เร็วๆ นี้

เพื่อที่จะขายหุ้นในราคาที่ดีให้กับนักลงทุน

บริษัทต้องการทำงานบางโครงการเพื่อเพิ่มทุนก่อนการเสนอขายหุ้น

เนื่องจากมีทรัพยากรที่จำกัดจึงสามารถดำเนินการให้เสร็จสิ้นในโครงการที่แตกต่างกันได้มากที่สุด\texttt{K}โครงการก่อนการเสนอขายหุ้น

ขอให้ท่านช่วยบริษัทออกแบบวิธีที่ดีที่สุดในการเพิ่มทุนทั้งหมดหลังจากเสร็จสิ้นโครงการที่แตกต่างกันมากที่สุด\texttt{K}โครงการ



คุณได้รับ\texttt{N}โครงการ

โดยโครงการที่\texttt{i}มีกำไรสุทธิ\texttt{profits{[}i{]}}และจำเป็นต้องมีเงินทุนขั้นต่ำ\texttt{capital{[}i{]}}เพื่อเริ่มต้นโครงการ



เริ่มแรกคุณมีเงินทุน\texttt{W}เมื่อคุณเสร็จสิ้นแต่ละโครงการ

คุณจะได้รับผลกำไรสุทธิข้างต้นและถูกบวกเพิ่มเข้าไปกับเงินทุนทั้งหมดของคุณ



ให้เลือกรายชื่อโครงการที่แตกต่างกันมากที่สุด\texttt{K}โครงการจากโครงการที่กำหนดเพื่อเพิ่มทุนสุดท้ายของคุณให้สูงสุด

และ\textbf{\ul{ตอบ}}เป็นค่าเงินทุนสูงสุดสุดท้ายนั้น (ทั้งนี้

คำตอบจะไม่เกินเลขจำนวนเต็มแบบ 32 บิต)



\hfill\break



{ตัวอย่างที่ 1}



\begin{quote}

\textbf{Input:}



2 0



3



1 2 3



0 1 1



\textbf{Output:}



4\\



{คำอธิบาย}



{\texttt{K\ =\ 2},\texttt{W\ =\ 0},\texttt{profits\ =\ {[}1,\ 2,\ 3{]}},\texttt{capital\ =\ {[}0,\ 1,\ 1{]}}}\\



เนื่องจากเงินทุนเริ่มต้นของคุณคือ 0 คุณจึงสามารถเริ่มโครงการที่ 0 ได้เท่านั้น



หลังจากเสร็จสิ้น คุณจะได้รับกำไร 1 และเงินทุนของคุณจะกลายเป็น 1



ด้วยทุน 1 คุณสามารถเริ่มโครงการที่จัดทำดัชนี 1 หรือโครงการที่จัดทำดัชนี 2 ได้



เนื่องจากคุณสามารถเลือกได้สูงสุด 2 โครงการ คุณจึงต้องดำเนินโครงการที่จัดทำดัชนี 2

ให้เสร็จสิ้นเพื่อรับทุนสูงสุด ดังนั้น

ให้ตอบเงินทุนสูงสุดสุดท้ายซึ่งคือ\texttt{0\ +\ 1\ +\ 3\ =\ 4}

\end{quote}



\subsection*{Input}
บรรทัดที่ 1 เป็นค่า\texttt{K}และ\texttt{W}เว้นด้วยช่องว่าง 1 ช่อง



{บรรทัดที่ 2

เป็นค่า\texttt{n}ขนาดของอาร์เรย์\hspace{0pt}{\texttt{capital}และ\texttt{profits}มีขนาดเท่ากัน

เท่ากับ\texttt{n}}}\\



{{บรรทัดที่ 3 เป็นค่าของอาร์เรย์\hspace{0pt}{\texttt{profits}}}\\

}



{{{บรรทัดที่ 4{เป็นค่าของอาร์เรย์ \hspace{0pt}{\texttt{capital}}}{}}\\

}}



\subsection*{Output}
{เงินทุนสูงสุดสุดที่ได้รับ}\\



\end{document}
